\scanpage{8}
\mspage{( 3 )}

\textsc{Corollaire.} --- \textit{Si $\mathrm A$ est une localité sur le corps $\mathrm K$, telle que $\mathrm K_a$ et $\mathrm L_a$ soient séparables sur $\mathrm K$, il y a équivalence entre :}

(i) \textit{$\mathrm D_{\mathrm K}(\mathrm A)$ est libre;}

(ii) $\mathrm A$ est régulier.

\textsc{Proposition 1.} --- \textit{Si $\mathrm A$ domine un corps $\mathrm K$, tel que $\mathrm D_{\mathrm K}(\mathrm A)$ soit libre et de rang fini, si $\mathrm K_b$ (resp. $\mathrm L_a$) est séparable sur $\mathrm K_a$ (resp. $\mathrm L_b$) et si}
\[
\operatorname{rg}\mathfrak n-\operatorname{rg}\mathfrak m=\{\mathrm K_b:\mathrm K_a\}-\{\mathrm L_b:\mathrm L_a\},
\]
\textit{$\mathrm D_{\mathrm A}(\mathrm B)$ est un $\mathrm B$-module de dimension homologique finie $\leq1$.}

\textsc{Corollaire.} --- \textit{Si $\mathrm A$ est un anneau local régulier contenant un corps $\mathrm K$, tel que le couple $(\mathrm K,\mathrm A)$ jouisse des propriétés énoncées dans le corollaire ci-dessus, $\mathrm D_{\mathrm A}(\mathrm B)$ est de dimension homologique $\leq1$ sous les hypothèses de la proposition 1.}

\textsc{Proposition 2.} --- \textit{Soit $\mathrm A$ un anneau local nœthérien contenant un corps $\mathrm K$ tel que $\{\mathrm L_a:\mathrm K\}$ soit fini et que $\mathrm L_a$ soit séparable sur $\mathrm K$. Si $\mathrm D_{\mathrm K}(\mathrm A)$ est libre, de rang $\dim\mathrm A+\{\mathrm L_a:\mathrm K\}$, $\mathrm A$ est local régulier.}

Notons, pour conclure, que dans de nombreuses questions (\textit{cf.} géométrie formelle) il vaut mieux séparer $\mathrm D_{\mathrm A}(\mathrm B)$ en considérant
\[
\mathrm D'_{\mathrm A}(\mathrm B)=\frac{\mathrm D_{\mathrm A}(\mathrm B)}{\displaystyle\bigcap_{n\geq0}\mathfrak n^n\mathrm D_{\mathrm A}(\mathrm B)};
\]
la proposition 2 est encore vraie, lorsqu’on y substitue $\mathrm D'_{\mathrm K}(\mathrm A)$ à $\mathrm D_{\mathrm K}(\mathrm A)$.

(\textsuperscript{*}) Séance du 14 janvier 1963.

(\textsuperscript{1}) \textsc{Berger} et \textsc{Kunz}, \textit{Math. Z.}, \textbf{77}, 1961, p. 314-338.

(\textsuperscript{2}) \textsc{Bourbaki}, \textit{Algèbre}, chap. V, \S~8, ex. 1.

(\textsuperscript{3}) \textsc{Bourbaki}, \textit{Algèbre comm.}, chap. I, \S~2, ex. 17.

(\textsuperscript{4}) \textsc{Godement}, \textit{Séminaire Cartan-Chevalley}, 1955-1956, exposés 17 et 18.

(\textsuperscript{5}) \textsc{Nakai}, \textit{J. Math. Soc. Japan}, 1961.

(\textsuperscript{6}) \textsc{Nakai} et \textsc{Suzuki}, \textit{J. Sc. Hiroshima Univ.}, série A, \textbf{24}, n\textsuperscript{o} 3, décembre 1960.

\hfill(\textit{École Normale Supérieure,} 45, \textit{rue d’Ulm, Paris,} 5\textsuperscript{e}.)
% archive sheet number: 7
